I tested the coupled W2L walker the way the perturbation literature tests neuromechanical models: the tuned conductances stay frozen, the perturbation is applied to the body or to the network, and recovery is judged from the gait itself (Figure~\ref{fig:w2l_perturbation}). Under a harness carrying \qty{25}{\percent} of body weight, the walker withstood every segment-mass change tested (thigh masses and distal-leg masses each scaled by $+10$, $-10$, $+20$, and $-20$ percent, and the whole leg by $\pm 10$ percent): ten of ten conditions, without a fall and with forward progression in each. An inhibitory tonic stimulus applied to the rhythm-generator layer for one second perturbed the rhythm without a fall, and stepping persisted through and after the stimulus window; a \qty{40}{\newton} fore-aft push at the pelvis likewise produced a transient excursion and recovery. Removing individual pathway classes then located the load-bearing connections. Deleting the Ia reciprocal chain, the Renshaw cells, the heel-contact cascade, or the crossed heel-to-contralateral-flexor edge left alternating stepping intact, whereas removing the commissural coupling reduced the gait to two strides and removing the Ib load feedback stopped stepping entirely. The same conductance set walks with the harness weight support re-tuned down a ladder from 70 to 10 percent of body weight, with a fresh-process verification at every step. These tests concern supported progression under the harness; unsupported walking is not established.

\begin{figure}[p]
\centering
\includegraphics[width=\textwidth,height=0.52\textheight,keepaspectratio]{w2l_perturbation_panels.pdf}
\caption[W2L perturbation and pathway-deletion panels]{Perturbation tests of the corrected W2L walker in MuJoCo with all conductances frozen (harness at \qty{25}{\percent} weight support). (A--C) Left and right knee flexion (solid and dashed; heel contact shaded). (B, C) A one-second inhibitory tonic stimulus to the rhythm-generator layer (yellow band), mild and strong; stepping persists through and after the window. (D, E) Fore and aft \qty{40}{\newton} pelvis pushes at $t=\qty{8}{\second}$ (dotted line): fore-aft center-of-mass position (solid) and pelvis tilt (dashed). (F) Strides per leg when one pathway class is removed: most removals are compensated, while commissural coupling and Ib load feedback are load-bearing.}
\label{fig:w2l_perturbation}
\end{figure}

